\documentclass{article}
\usepackage{graphicx} % Required for inserting images

\usepackage{tikz-cd}


%math symbols
\usepackage{graphicx, amsmath, amssymb, longtable, mathtools, amsthm, polynom, mathrsfs, cancel}
%\usepackage{enumitem}

%\newcommand{\tabitem}{~~\llap{\textbullet}~~}
\newcommand{\tabitem}{\textbullet~}
\usepackage{url}
\usepackage{breakcites}
\usepackage{comment}
\usepackage[table]{xcolor}
\usepackage{booktabs} % Optional: for better looking horizontal lines

%Bibtex
\usepackage[backend=biber,style=alphabetic]{biblatex}
\addbibresource{references.bib}
%\usepackage[breaklinks]{hyperref}
%\appto\bibfont{\setlength{\emergencystretch}{.1em}}


%Theorem Like Environments

%\newtheorem{theorem}{Theorem} 
%\numberwithin{theorem}{section}




%\newtheorem{definition}{Definition}
%\numberwithin{definition}{section}

\newtheorem{assumption}{Assumption}
\numberwithin{assumption}{section}

%\newtheorem{proposition}{Proposition}
%\numberwithin{proposition}{section}

%\newtheorem{remark}{Remark}
%\numberwithin{remark}{section}

%\newtheorem{corollary}{Corollary}
%\numberwithin{corollary}{section}

%\newtheorem{conjecture}{Conjecture}
%\numberwithin{conjecture}{section}

%\newtheorem{example}{Example}
%\numberwithin{example}{section}


%\usepackage{comment}
%\renewenvironment{comment}{}{}
\usepackage{algorithm}
\usepackage{algpseudocode} 
%\usepackage{algorithm2e}

%Tabular packages
\usepackage{multicol, multirow, array, tabularx, caption, float}
\newcolumntype{C}[1]{>{\centering\arraybackslash}p{#1}}
%Tikz package
\usepackage{tikz}
\usetikzlibrary{positioning}
\usepackage{subcaption}

%set margins



%Common sets and symbols:
\newcommand{\R}{\mathbb{R}}
\newcommand{\Z}{\mathbb{Z}}
\newcommand{\N}{\mathbb{N}}
\newcommand{\Q}{\mathbb{Q}}
\newcommand{\C}{\mathbb{C}}
\newcommand{\F}{\mathbb{F}}
\newcommand{\bigS}{\mathcal{S}}

\newcommand{\Open}{\mathcal{O}}
\newcommand{\E}{\mathbb{E}}
\newcommand{\I}{\mathbb{I}}
\newcommand{\Var}{\mathbf{Var}}
\newcommand{\bset}[1]{\left \lbrace #1 \right \rbrace}
\newcommand{\aset}[1]{\left \langle #1 \right \rangle}
\newcommand{\mexists}{\; \exists \;}
\newcommand{\mforall}{\; \forall \;}
\newcommand{\cunion}[2]{\bigcup_{#1 = #2}^\infty}
\newcommand{\cinter}[2]{\bigcap_{#1 = #2}^\infty}
\newcommand{\csum}[2]{\sum_{#1 = #2}^\infty}
\newcommand{\dist}{\mathrm{dist}}
\newcommand{\Span}[1]{\mathrm{Span}\bset{#1}}
\newcommand{\id}{\mathrm{id}}
\newcommand{\im}{\mathrm{im}}
\newcommand{\D}{\displaystyle}
\newcommand{\Cov}{\mathrm{Cov}} 
\newcommand{\intp}[1] {\D \int_\Omega {#1 } \; d \mathbb P}
\newcommand{\HS}{\mathrm{HS}}

\newcommand{\NT}[2]{N_\mathbf{#1}^{\text{#2}}}
\newcommand{\T}[2]{\mathcal T_\mathbf{#1}^{\text{#2}}}

\newcommand{\bA}{\mathbf A}
\newcommand{\bB}{\mathbf B}
\newcommand{\CH}{\mathcal H}
\newcommand{\Mres}{M_\text{res}}
%\renewenvironment{comment}{}{}

\newcommand{\globalLGheight}{0.23\textheight}
\newcommand{\globalLGwidth}{0.87\textwidth}
\title{Option Pricing For Weather-Based Derivatives}
\author{Trajan Murphy}
\date{\today}

%\includeonly{Experimental Details/Experimental Results.tex}


\begin{document}
	\maketitle
	
	Weather-sensitive industries, such as air travel, agriculture, and tourism often incorporate weather-based derivatives into their models for option pricing. We consider a simple model, based on temperature derivatives, by following the thesis work of Gilbert, 2024 \cite{gibert2024monte}. We will first begin by modeling the temperature $T_t$ at time $t$, via an SPDE known as a \textit{mean-reverting OU Process}, or simply OU process. The OU process is given by
	
		\begin{equation} \label{Mean-reverting OU-Process}
		dT_t = \kappa(\overline T_t - T_t) + \sigma_t dW_t
	\end{equation}
	
	where 
	
	\begin{itemize}
		\item $T_t$ = Temperature at time $t$
		\item $\overline T_t$ = mean temperature at time $t$
		\item $\kappa$ = ``constant"
		\item $\sigma_t$ = volatility at time $t$
		\item $W_t$ = Wiener process
	\end{itemize}
	
	An explicit solution is given by the stochastic integral: 
	
	\begin{equation} \label{Explicit Solution}
		T_t = \overline T_t + \exp \left( -\int_0^t \kappa du \right) \int_0^t \exp \left( \int_0^t \kappa du \right) \sigma_t dW_t.
	\end{equation}
	
	$\overline T_t$ is typically estimated from historical data. In our simulation, we fit the daily dry bulb temperature recorded at the Boston Logan Airport in the year 2012, sourced from \cite{NCEI2026} to a quadratic model. Once we obtain $T_t$ we compute two different underlying indices based on the time of year: the \textit{Heating Degree Days} and \textit{Cooling Degree Days}, given respectively by 
	
	\begin{equation}
		\begin{split}
			CDD_n = (T_n - 65)^+
			\\
			HDD_n = (65 - T_n)^+
		\end{split}
	\end{equation}\\
	
	where $(\cdot)^+ = \max{\cdot, 0}$. Here $n$ is the \textit{discretized} time index. The HDD is used during the heating season (typically defined to be May 15-Oct 15), and the CDD is used during the cooling season (typically defined to be Nov 15-Mar 15). Once we have these sequences, the value of a call (respectively put) option at day $N$ is given by 
	
	\begin{equation} \label{Option Prices}
		\begin{split}
			\text{Call} &= \min \bset{\alpha \left( \left(\sum_{n=1}^N DD_n \right) - K_c \right)^+, C}
			\\
			\text{Put} &= \min \bset{\alpha \left(K_c - \left(\sum_{n=1}^N \right)   \right)^+, C}
		\end{split}
	\end{equation}
	
	where Let $K_c$ is the strike price, $\alpha$ is the payout rate, and $C$ is the cap price. Typical values of these constants are $2,500 \leq \alpha \leq 5,000$ and $500,000 \leq C \leq 1,000,000$. In our simulations we fix $\alpha = 2,500$, $C = 500,000$, and $K_C = 550$. 
	
	To produce our Monte-Carlo simulations of the daily call/put option prices, we do the following. We estimate the daily average temperature using trapezoidal integration on the hour-by-hour values. We fit the daily average temperature to a quadratic model. We compute the ACF, PACF, and QQ-plot of residuals to ensure normality and independence. We estimate $\sigma_t$ by using a monthly average and linearly interpolate to get daily values. We then simulate multiple $T_t$ series with $\kappa \equiv 1$ according to \eqref{Explicit Solution}. Finally, we compute call and put options based on modeled temperature according to \eqref{Option Prices}.
	
	\printbibliography
\end{document}